\chapter{Overview of the Lattice Approach}

LQCD calculations are a non-perturbative implementation of field theory
using the Feynman path integral approach. The calculations proceed exactly as if the field theory was being solved analytically had we the ability
to do the calculations. The starting point is the partition function in Euclidean space-time
\begin{equation}
  Z = \int \mathcal{D}A_\mu\, \mathcal{D}\bar{\psi}\, \mathcal{D}\psi\, e^{-S}
  \label{eq:partition}
\end{equation}
where $S$ is the QCD action
\begin{equation}
  S = \int d^4x \left( \frac{1}{4} F_{\mu\nu}F^{\mu\nu} - \bar{\psi} M \psi \right).
  \label{eq:qcd_action}
\end{equation}
and $M$ is the Dirac operator. The fermions are represented by Grassmann
variables $\psi$ and $\bar{\psi}$. These can be integrated out exactly with
the result
\begin{equation}
  Z = \int \mathcal{D}A_\mu\, \det M\, e^{\int d^4x \left(-\frac{1}{4} F_{\mu\nu}F^{\mu\nu}\right)}.
  \label{eq:partition_det}
\end{equation}
The fermionic contribution is now contained in the highly non-local term
$\det M$, and the partition function is an integral over only background
gauge configurations. One can write the action, after integration over the
fermions, as $S = S_\mathrm{gauge} + S_\mathrm{quarks} = \int d^4x ( \frac{1}{4} F_{\mu\nu}F^{\mu\nu} - i \sum_i \log(\mathrm{Det}\,M_i) )$
where the sum is over the quark flavors, distinguished by the value of the
bare quark mass.

It is expedient to define the ``quenched'' approximation (QQCD) from
the very start due to the central role it has played in simulations. It consists of setting $\mathrm{Det}\,M = \mathrm{constant}$ which corresponds to removing vacuum
polarization effects from the QCD vacuum. Details and some consequences
of this approximation are discussed in Section 16.8.

Results for physical observables are obtained by calculating expectation
values
\begin{equation}
  \langle O \rangle = \frac{1}{Z} \int \mathcal{D}A_\mu\, O\, e^{-S}.
  \label{eq:expectation}
\end{equation}
where $O$ is any given combination of operators expressed in terms of timeordered products of gauge and quark fields. (Expectation values in the
path integral approach correspond to time-ordered correlation functions.)
The quarks fields in $O$ are, in practice, re-expressed in terms of quark
propagators using Wick's theorem for contracting fields. In this way all
dependence on quarks as dynamical fields is removed. The basic building
block for fermionic quantities is the Feynman propagator,
\begin{equation}
  S_F(y, j, b; x, i, a) = \left( M^{-1} \right)^{y,j,b}_{x,i,a},
  \label{eq:feynman_propagator}
\end{equation}
where $M^{-1}$ is the inverse of the Dirac operator calculated on a given background field. A given element of this matrix $\left( M^{-1} \right)^{y,j,b}_{x,i,a}$ is the amplitude for
the propagation of a quark from site $x$ with spin-color $i, a$ to site-spin-color
$y, j, b$.

So far all of the above is standard field theory. The problem we face
in QCD is how to actually calculate these expectation values and how to
extract physical observables from these. I will illustrate the second part
first by using as an example the mass and decay constant of the pion.

Consider the 2-point correlation function, $\langle 0| T \left[ \sum_x O_f(x, t) O_i(0, 0) \right] |0\rangle$
with $t > 0$, where the operators $O$ are chosen to be the fourth component of
the axial current $O_f = O_i = A_4 = \bar{\psi}\gamma_4\gamma_5\psi$ as these have a large coupling to
the pion. This 2-point correlation function gives the amplitude for creating
a state with the quantum numbers of the pion out of the vacuum at spacetime point $0$ by the ``source'' operator $O_i$; the evolution of this state to
the point $(x, t)$ via the QCD Hamiltonian (or more precisely the transfer
matrix to be discussed in Section 5); and finally the annihilation by the
``sink'' operator $O_f$ at $(x, t)$. The rules of quantum mechanics tell us that
$O_i$ will create a state that is a linear combination of all possible eigenstates
of the Hamiltonian that have the same quantum numbers as the pion, i.e.\
the pion, radial excitations of the pion, three pions in $J = 0$ state, \ldots.
The second rule is that on propagating for Euclidean time $t$, a given eigenstate
with energy $E$ picks up a weight $e^{-Et}$. Thus, the 2-point function can be
written in terms of a sum over all possible intermediate states
\begin{equation}
  \langle 0| \sum_x O_f(x, t) O_i(0) |0\rangle
  = \sum_n \frac{\langle 0|O_f|n\rangle\langle n|O_i|0\rangle}{2E_n} e^{-E_n t}.
  \label{eq:2pt_spectral}
\end{equation}
To study the properties of the pion at rest we need to isolate this state from
the sum over $n$. To do this, the first simplification is to use the Fourier
projection $\sum_x$ as it restricts the sum over states to just zero-momentum
states, so $E_n \to M_n$. (Note that it is sufficient to make the Fourier projection over either $O_i$ or $O_f$.) The second step to isolate the pion, i.e.\
project in the energy, consists of a combination of two strategies. One,
make a clever choice of the operators $O$ to limit the sum over states to a
single state (the ideal choice is to set $O$ equal to the quantum mechanical
wave-functional of the pion), and two, examine the large $t$ behavior of the
2-point function where only the contribution of the lowest energy state that
couples to $O_i$ is significant due to the exponential damping. Then
\begin{equation}
  \langle 0| \sum_x O_f(x, t) O_i(0) |0\rangle
  \xrightarrow[t\to\infty]{}
  \frac{\langle 0|O_f|\pi\rangle\langle\pi|O_i|0\rangle}{2M_\pi} e^{-M_\pi t}.
  \label{eq:2pt_pion}
\end{equation}
The right hand side is now a function of the two quantities we want since
$\langle 0|A_4(p = 0)|\pi\rangle = M_\pi f_\pi$. In this way, the mass and the decay constant
are extracted from the rate of exponential fall-off in time and from the
amplitude.

Let me now illustrate how the left hand side of Eq.~\eqref{eq:2pt_pion} is expressed in
terms of the two basic quantities we control in the path integral --- the gauge
fields and the quark propagator. Using Wick contractions, the correlation
function can be written in terms of a product of two quark propagators $S_F$,
\begin{equation}
  \begin{aligned}
    -&\langle 0| \sum_x \bar{\psi}(x, t)\gamma_4\gamma_5\psi(x, t)\ \bar{\psi}(0, 0)\gamma_4\gamma_5\psi(0, 0)|0\rangle \\
    &\equiv \langle 0| \sum_x S_F(0; x, t)\gamma_4\gamma_5 S_F(x, t; 0)\gamma_4\gamma_5|0\rangle.
  \end{aligned}
  \label{eq:wick_local}
\end{equation}

\begin{figure}[htbp]
  \centering
  \includegraphics[width=0.8\textwidth]{images/fig4.png}
  \caption{A schematic of the pion 2-point correlation function for (A) local and (B) nonlocal interpolating operators.}
  \label{fig:4}
\end{figure}

This correlation function is shown in Fig.~\ref{fig:4}A. To illustrate this Wick contraction procedure further consider using gauge
invariant non-local operators, for example take $O = \bar{\psi}(x, t)\gamma_4\gamma_5 \left( P e^{i g \int_x^y dz\, A_\mu(z)} \right)\psi(y, t)$ where $P$
stands for path-ordered. After Wick contraction the correlation function
reads
\begin{equation}
  \langle 0| \sum_x S_F(0; x, t)\gamma_4\gamma_5 \left( P e^{i g \int_z^x A_\mu} \right) S_F(z, t; y, 0)\gamma_4\gamma_5 \left( P e^{i g \int_0^y A_\mu} \right)|0\rangle.
  \label{eq:wick_extended}
\end{equation}
and involves both the gauge fields and quark propagators. This correlation function would have the same long $t$ behavior as shown in Eq.~\eqref{eq:2pt_spectral},
however, the amplitude will be different and consequently its relation to
$f_\pi$ will no longer be simple. The idea behind strategies for improving
the projection of $O$ on to the pion is to construct a suitable combination of such operators that approximates the pion wave-function. For example, if $\phi(x)$ is the pion wave-functional,
then the ``smeared'' operator $O = \bar{\psi}(0)\gamma_4\gamma_5 \int d^3x\, \phi(x)\psi(x)$ is the ideal choice.

Finally, it is important to note that to get the 2-point function corresponding to the propagation of the physical pion we have to average the
correlation function over gauge configurations as defined in Eq.~\eqref{eq:expectation}, i.e.\ do
the functional integral.

An overview of how such calculations of expectation values can be done
consists of the following steps.
\begin{itemize}
  \item \textbf{Defining a finite dimensional system:} The Yang-Mills action for gauge
    fields and the Dirac operator for fermions have to be transcribed onto a discrete space-time lattice in such a way as to preserve all the key properties
    of QCD --- gauge invariance, chiral symmetry, topology, and a one-to-one
    relation between continuum and lattice fields. This step is the most difficult, and even today we do not have a really satisfactory lattice formulation
    that is chirally symmetric in the $m_q = 0$ limit and preserves a one-to-one
    relation between continuum and lattice fields, i.e.\ no doublers. In fact, the
    Nielson-Ninomiya theorem states that for a translationally invariant, local,
    hermitian formulation of the lattice theory one cannot simultaneously have
    chiral symmetry and no doublers [14]. One important consequence of this
    theorem is that, in spite of tremendous effort, there is no practical formulation of chiral fermions on the lattice. For a review of the problems and
    attempts to solve them see [15--17] and [18].
  \item \textbf{Discretization Errors:} This problem is encountered when approximating
    derivatives in the action by finite differences,
    \begin{equation}
      \partial_x \phi(x) \to \Delta_x \phi(x)
      = \frac{1}{2a}\left( \phi(x + a) - \phi(x - a) \right)
      = \frac{1}{2a}\left( e^{a\partial_x} - e^{-a\partial_x} \right)\phi(x).
      \label{eq:derivative}
    \end{equation}
    As is obvious this introduces discretization errors proportional to the square
    of the lattice spacing $a$. These errors can be reduced by either using higher
    order difference schemes with coefficients adjusted to take into account
    effects of renormalization, or equivalently, by adding appropriate combinations of irrelevant operators to the action that cancel the errors order by order in $a$. The various approaches to improving the gauge and Dirac actions
    are discussed in Sections 12 and 13. In my lectures I will concentrate on
    the three most frequently used discretizations of the Dirac action --- Wilson
    [19], Sheikholeslami-Wohlert (clover) [20], and staggered [21], which have
    errors of $O(a)$, $O(\alpha_s a) - O(a^2)$ depending on the value of the coefficient
    of the clover term, and $O(a^2)$ respectively. Reducing the discretization errors by improving the action and operators increases the reliability of the
    results, and for dynamical simulations may even prove necessary due to
    limitations of computer power. At the same time it is important to note
    that there exists a well defined procedure for obtaining continuum results
    with even the simplest formulation, Wilson's original gauge and fermion
    action [19]. In this case, however, it is necessary to quantify and remove
    the discretization errors by a reliable extrapolation to $a = 0$.
  \item \textbf{Generating background gauge fields:} The Euclidean action $S = \int d^4x ( \frac{1}{4} F_{\mu\nu}F^{\mu\nu} - \mathrm{Tr}\log M )$ for QCD at zero chemical potential is real
    and bounded from below. Thus $e^{-S}$ in the path integral is analogous to
    the Boltzmann factor in the partition function for statistical mechanics
    systems, i.e.\ it can be regarded as a probability weight for generating configurations. Since $S$ is an extensive quantity the range of $S$ is enormous,
    however, the distribution is very highly peaked about configurations that
    minimize the free energy. Configurations close to the peak dominate the
    functional integral. Such important configurations can be generated by
    setting up a Markov chain in exact analogy to, say, simulations of the Ising
    model. They are called ``importance sampled'' as they are generated with
    probability of occurrence given by the weight $e^{-S}$. For a discussion of the
    methods used to update the configurations see [1] or the lectures by Creutz
    and Sokal in [2].
  \item \textbf{Calculation of the quark propagator on a given background gauge field:} For
    a given background gauge configuration, the Feynman quark propagator $S_F$
    is a matrix labeled by three indices --- site, spin, and color. A given element
    of this matrix gives the amplitude for the propagation of a quark with
    some spin, color, and space-time point to another space-time point, spin,
    and color. Operationally, $S_F$ is simply the inverse of the Dirac operator
    $M$. Once space-time is made discrete and finite, the Dirac matrix is also
    finite and its inverse can be calculated numerically. Inverting the Dirac
    operator is the key computational step in LQCD and consumes over 90\%
    of the CPU cycles. To construct the correlation functions requires $S_F$ only
    from a given source point to all other sites of the lattice. This corresponds
    to 12 columns (or rows) of $M^{-1}$, one each for the 12 spin and color degrees
    of freedom.
  \item \textbf{Correlation functions:} The correlation functions are expressed as path
    ordered products of quark propagators and gauge fields using Wick contraction of the operator fields. There are two kinds of correlation functions that
    are calculated --- gauge invariant correlation functions as discussed above
    for the case of the pion, or those in a fixed gauge.
  \item \textbf{Expectation values:} The expectation value is simply an average of the
    correlation function evaluated over the set of ``importance sampled'' configurations. Even on a finite lattice the set of background gauge configurations is infinite as the variables are continuous valued. Thus, while it
    is possible to calculate the correlation functions on specified background
    gauge configurations, doing the functional integral exactly is not feasible.
    It is, therefore, done numerically using monte carlo methods.
\end{itemize}

My recommendation for the simplest way to understand the computational aspects of LQCD calculations is to first gain familiarity with the
numerical treatment of any simple statistical mechanics system, for example the Ising model. The differences are: (i) the degrees of freedom in
LQCD are much more complicated --- SU(3) link matrices rather than Ising
spins, and quark propagators on these background configurations evaluated
by inverting the Dirac operator; (ii) The action involves the highly nonlocal
term Ln Det $M$ which makes the update of the gauge configurations very
expensive; and (iii) the correlation functions are not simple products of spin
variables like the specific heat or magnetic susceptibility, but complicated
functions of the link variables and quark propagators.

The subtleties arising due to the fact that LQCD is a renormalizable
quantum field theory and not a classical statistical mechanics system come
into play when analyzing the behavior of the correlation functions as a
function of the lattice spacing $a$, and in the quantum corrections that renormalize the input parameters (quark and gluon fields and their masses) and
the composite operators used in the correlation functions. At first glance
it might seem that we have introduced an additional parameter, the lattice
spacing $a$, in LQCD, however, as will be shown in Section 14, the coupling
$g$ and the cutoff $a$ are not independent quantities but are related by the
renormalization group as
\begin{equation}
  \Lambda_\mathrm{QCD} = \lim_{a\to 0} \frac{1}{a} e^{-1/(2\beta_0 g^2(a))} \left( \beta_0 g^2(a) \right)^{-\beta_1/(2\beta_0^2)},
  \label{eq:lambdacd}
\end{equation}
where $\Lambda_\mathrm{QCD}$ is the non-perturbative scale of QCD, and $\beta_0$ and $\beta_1$ are the
first two, scheme independent, coefficients of the $\beta$-function. In statistical
mechanics systems, the lattice spacing $a$ is a physical quantity --- the intermolecular separation. In a quantum field theory, it is simply the ultraviolet
regulator that must eventually be taken to zero keeping physical quantities,
like the renormalized couplings ($\Lambda_\mathrm{QCD}$ and quark masses), spectrum, etc,
fixed.

The reason that lattice results are not exact is because in numerical
simulations we have to make a number of approximations as discussed in
Section 16. The size of the associated uncertainties is dictated by the
computer power at hand. They are being improved steadily with computer
technology, better numerical algorithms, and better theoretical understanding. I shall return to the issue of errors after discussing the formulation of
LQCD.
